\section{How to run}
\label{s:howto}

\begin{lstlisting}
# input: a file list, one NanoAOD path per line (see Sec. algorithm: input branches)
./run.sh p1 8                  # pass 1, all events, no correction, 8 jobs -> rootfiles/GenSeed_p1.root
root -l -b -q 'drawGenSeed.C+("rootfiles/GenSeed_p1.root","_v1")'   # -> text/jec_v1.txt, plots/, doc/{plots,tables}.tex
./run.sh p2 8                  # pass 2 with text/jec_v1.txt -> rootfiles/GenSeed_v1.root
root -l -b -q 'drawGenSeed.C+("rootfiles/GenSeed_v1.root","_v1")'
cd doc && pdflatex genseed && pdflatex genseed
\end{lstlisting}

\noindent A single job on a range of entries, or a test:
\begin{lstlisting}
root -l -b -q 'runGenSeed.C("files.txt","_test",0,1,200)'          # 200 events
root -l -b -q 'runGenSeed.C("files.txt","_v1",3,8,0,"text/jec_v1.txt")'   # job 3 of 8
\end{lstlisting}

\noindent The clustering backend is chosen at build time in \texttt{runGenSeed.C}:
\texttt{tiledjet.h} unless FastJet is found \emph{and} the benchmark constant
\texttt{tiledjet::kAtLeastAsFastAsFastJet} is false. \texttt{bench\_tiledjet.C}
reproduces the benchmark on any machine. Every knob is a \texttt{SetOpt} call
in the driver and every value used is written to \texttt{hist/hopts}, with the
PUPPI version and the correction table as text next to it
(\texttt{hist/puppiVersion}, \texttt{hist/jecFile}), so a histogram file
always says how it was made.

\noindent What is in the repository. The input sample is not public, so the
outputs a reader needs are committed with the code: \texttt{plots/*.pdf},
the generated \texttt{doc/plots.tex} and \texttt{doc/tables.tex}, the
compiled \texttt{doc/genseed.pdf}, and the correction table
\texttt{text/jec\_v1.txt} that pass 2 reads. Anyone can read the document, or
rebuild it with \texttt{cd doc \&\& pdflatex genseed}, without the sample.
The histogram and tuple files (\texttt{rootfiles/}), everything else in
\texttt{text/}, the ACLiC build products and the \LaTeX{} auxiliary files are
not committed. After a new production, rerun \texttt{drawGenSeed.C} and
\texttt{pdflatex} and commit the plots, the two generated \texttt{.tex} files
and the PDF together, so that they never describe different files.

\section{Towards JMENANO}
\label{s:jmenano}

The algorithm needs, per event: the PF candidates with $\pt$, $\eta$, $\phi$,
mass, pdgId, charge, $dz$ to the primary vertex, the vertex key of each charged
candidate (\texttt{vertexRef}) with the fit-quality flags, and the generated
particles of \emph{every} pileup interaction with their interaction $z$. The
first set is what \texttt{PFCandV2} in the custom NanoAOD used here carries
and what JMENANO's PF candidate table carries or can be extended to carry
(\texttt{vertexRef} and \texttt{dz} are the essential additions). The second
set is only available in a production that keeps the pileup generator record,
as this sample does; for a standard production with the signal generated
separately from the pileup, the linker runs on the signal's generated particles
alone, the GenJetSeed is defined as here with everything unlinked counted as
pileup, and the RecoJetSeed is defined for the signal jets. Nothing in
\texttt{genlink.h} or \texttt{GenSeed.C} depends on the pileup record beyond
the ownership of a vertex, which then reduces to the signal vertex.
